% =====================================================================
%  2bat-ud6-15.tex  —  SESSIÓ 15: Taller de preparació de la prova
%  (sense preàmbul: s'inclou des de main.tex)
% =====================================================================
\horatitol{Sessió 15 — Taller de preparació de la prova}%
{Conèixer l'estructura de la prova i la rúbrica, assajar un simulacre per blocs i fer un pla d'estudi personal.}

\subsection{Idea clau}

Ens preparem per a la prova \textbf{sense sorpreses}. La prova de la sessió següent —la que
tanca el curs i l'etapa— té \textbf{quatre blocs}, un per cada criteri (C1--C4). En aquesta
unitat, més que calcular, es valora \textbf{interpretar i expressar conclusions amb
prudència}. Després assajarem amb un \textbf{simulacre}.

\begin{idea}
\textbf{Estructura de la prova} (un bloc per criteri):
\begin{enumerate}[leftmargin=1.6em, itemsep=2pt]
  \item \textbf{Bloc 1 (C1) — Mostreig:} valorar la representativitat d'una mostra i
        detectar-ne un biaix.
  \item \textbf{Bloc 2 (C2) — Intervals:} construir i interpretar l'interval de confiança
        d'un resultat.
  \item \textbf{Bloc 3 (C3) — Significació:} valorar si una diferència entre dos grups és
        real o casual.
  \item \textbf{Bloc 4 (C4) — Lectura crítica:} analitzar un titular o fragment d'informe i
        detectar-ne els paranys.
\end{enumerate}
\textbf{Es valora sobretot:} la correcció de les \textbf{interpretacions}, la detecció de
\textbf{biaixos i paranys} i l'expressió de conclusions \textbf{sense sobreinterpretar}.
\end{idea}

\subsection{Rúbrica simplificada}

\noindent Així es corregeix cada bloc (nivells d'assoliment):
\begin{center}
\renewcommand{\arraystretch}{1.3}
\begin{tabular}{p{2.4cm} p{10.2cm}}
\toprule
\textbf{Nivell} & \textbf{Què demostra} \\
\midrule
\textbf{AE} (excel·lent) & Interpreta correctament i expressa la conclusió amb la prudència
adequada. \\
\textbf{AN} (notable) & Interpreta bé amb algun matís imprecís; la conclusió és
essencialment prudent. \\
\textbf{AS} (suficient) & Identifica la idea però amb errors d'interpretació o poca
prudència. \\
\textbf{NA} (no assolit) & Sobreinterpreta les dades o no detecta el biaix / el parany. \\
\bottomrule
\end{tabular}
\end{center}

\subsection{Simulacre — un ítem de cada bloc}

\begin{exemple}[com es resol un ítem ben fet]
\textbf{Ítem tipus (Bloc 2):} una enquesta dona «$46\%\pm 3\%$». Un \textbf{ítem complet}
construeix l'interval \emph{i} l'interpreta amb prudència: «l'interval és $[43\%,49\%]$; és
molt probable que el valor real de la població hi estigui, no que sigui exactament el
$46\%$». \emph{Expressar la conclusió amb prudència, i no com una certesa, distingeix un AE
d'un AS.}
\end{exemple}

\begin{exercicis}
\textbf{Resol el simulacre amb temps controlat. Un ítem per bloc:}
\begin{exlist}
  \item \textbf{(Bloc 1 — C1)} Una enquesta sobre satisfacció amb un servei es fa
        \textbf{només a qui hi continua anant}. Valora si la mostra és representativa i
        digues en quina direcció esbiaixa el resultat.
  \item \textbf{(Bloc 2 — C2)} Una enquesta dona «$52\%\pm 4\%$». Construeix l'interval de
        confiança i escriu què afirma correctament l'estudi.
  \item \textbf{(Bloc 3 — C3)} Dos programes donen un $55\%\pm 3\%$ i un $58\%\pm 3\%$
        d'èxit. Se superposen els intervals? La diferència és significativa?
  \item \textbf{(Bloc 4 — C4)} Titular: «\emph{Un estudi demostra que la població prefereix
        la mesura}», a partir d'una enquesta clicable en una web. Detecta els paranys.
\end{exlist}
\end{exercicis}

\subsection{Formulari-resum (suport)}

\begin{idea}
\textbf{Tingues-ho a mà mentre prepares la prova:}
\begin{itemize}[leftmargin=1.4em, itemsep=2pt]
  \item \textbf{Població i mostra:} població = tot el col·lectiu; mostra = la part
        mesurada. Ha de ser representativa (a l'atzar, sense deixar grups fora).
  \item \textbf{Interval de confiança:} resultat $\pm$ marge; «és molt probable que el
        valor real hi estigui». El nivell de confiança ($95\%$) no és «la gent».
  \item \textbf{Significació:} si els intervals \textbf{se superposen}, la diferència pot
        ser atzar; compta la mida de la diferència i de les mostres.
  \item \textbf{Lectura crítica:} d'on surten les dades?, mostra representativa?, expressa
        la incertesa?, correlació $\neq$ causa?, la diferència és significativa?
\end{itemize}
\end{idea}

\noindent\textbf{Pla d'estudi personal.} A partir de la graella d'autoavaluació (sessió 9) i
del simulacre, anota els \textbf{dos o tres aspectes} que has de repassar amb més atenció:
\begin{center}
\renewcommand{\arraystretch}{1.5}
\begin{tabular}{c p{10.5cm}}
\toprule
\textbf{\#} & \textbf{Què he de repassar (i com)} \\
\midrule
1 & \rule{9.8cm}{0.4pt} \\
2 & \rule{9.8cm}{0.4pt} \\
3 & \rule{9.8cm}{0.4pt} \\
\bottomrule
\end{tabular}
\end{center}

% ---------------------------------------------------------------------
\fullsolucions{Sessió 15 — simulacre}
\begin{solucions}
\begin{sollist}
  \item \textbf{(Bloc 1)} No és representativa: enquestar només qui continua anant-hi deixa
        fora qui va abandonar (potser per insatisfacció). \textbf{Sobreestima} la
        satisfacció (biaix de selecció).
  \item \textbf{(Bloc 2)} Interval $[48\%,\,56\%]$. Afirma correctament: és molt probable
        que la proporció real de tota la població estigui entre el $48\%$ i el $56\%$; no
        que sigui exactament el $52\%$.
  \item \textbf{(Bloc 3)} Intervals $[52\%,58\%]$ i $[55\%,61\%]$: \textbf{se superposen}
        (del $55\%$ al $58\%$). La diferència de $3$ punts \textbf{no} és significativa: pot
        ser atzar.
  \item \textbf{(Bloc 4)} Paranys: (1) mostra \textbf{autoseleccionada} (enquesta clicable,
        no representativa de la població); (2) «\textbf{demostra}» presenta com a certesa el
        que és, com a molt, una estimació poc fiable. El titular no es pot sostenir.
\end{sollist}
\end{solucions}
